\chapter*{ЗАКЛЮЧЕНИЕ}
\addcontentsline{toc}{chapter}{ЗАКЛЮЧЕНИЕ}

%% Согласно ГОСТ Р 7.0.11-2011:
%% 5.3.3 В заключении диссертации излагают итоги выполненного исследования, рекомендации, перспективы дальнейшей разработки темы.
%% 9.2.3 В заключении автореферата диссертации излагают итоги данного исследования, рекомендации и перспективы дальнейшей разработки темы.

В диссертационной работе рассмотрена задача адаптивного распределения вычислительных задач автономных мобильных роботов между периферийными вычислительными узлами. Данная задача возникает в роботизированных складских, производственных, терминальных и иных локальных AMR-системах, где группа роботов использует общую периферийную вычислительную инфраструктуру для выполнения задач восприятия, локализации, построения карты, обнаружения препятствий, перепланирования маршрута, инвентаризации и диагностики.

Проведённый анализ показал, что рассматриваемая задача не сводится только к статическому назначению задач на вычислительные узлы. Она имеет последовательный характер, поскольку текущее распределение задач влияет на будущие очереди, загрузку узлов, вероятность нарушения сроков обработки и доступность ресурсов для последующих критических задач. Дополнительную сложность создают частичная наблюдаемость, неоднородность задач, динамическая сетевая обстановка, ограниченность периферийных ресурсов и рост накладных расходов при увеличении числа роботов, задач и узлов.

В ходе исследования были получены следующие основные результаты.

%% Согласно ГОСТ Р 7.0.11-2011:
%% 5.3.3 В заключении диссертации излагают итоги выполненного исследования, рекомендации, перспективы дальнейшей разработки темы.
%% 9.2.3 В заключении автореферата диссертации излагают итоги данного исследования, рекомендации и перспективы дальнейшей разработки темы.

Во-первых, выполнен анализ предметной области автономных мобильных роботов с поддержкой периферийных вычислений. Показано, что существующие подходы к распределению ресурсов в периферийных системах развиваются по нескольким направлениям: детерминированная оптимизация, аукционные механизмы, обучение с подкреплением, многоагентное обучение и гибридные схемы. Детерминированные методы позволяют строго задавать ограничения и целевые функции, но требуют достаточно полной и устойчивой информации о состоянии системы. Аукционные механизмы позволяют учитывать дефицит ресурсов и внешние эффекты распределения, но в классическом виде ориентированы на отдельный раунд. Методы обучения с подкреплением позволяют учитывать долгосрочные последствия решений, но сами по себе не задают экономически интерпретируемого механизма распределения дефицитных ресурсов.

Во-вторых, разработана доменная и математическая модель AMR-системы с поддержкой периферийных вычислений. В модели определены множества автономных мобильных роботов, периферийных вычислительных узлов, вычислительных задач, граф сетевой связности, параметры задержек, очередей, отказов и ресурсных ограничений. Отдельно учтены классы робототехнических задач, различающиеся по критичности, срокам обработки, объёму данных и требованиям к ресурсам. Предложенная модель позволяет описывать не только вычислительную сложность задачи, но и её прикладную значимость для робота и текущей миссии.

В-третьих, построен VCG-подобный аукционный слой распределения вычислительных задач. В рамках одного раунда распределения задача назначения формализована как задача максимизации общественного благосостояния с учётом полезности выполнения задач, стоимости использования ресурсов, сетевых задержек, очередей, отказов и ограничений по срокам обработки. Платежи аукционного слоя интерпретируются как оценки внешнего эффекта, создаваемого участниками при использовании дефицитных ресурсов. При этом в работе явно зафиксированы границы применимости VCG-подобных свойств: эффективность, индивидуальная рациональность и стратегическая совместимость рассматриваются только для отдельного статического раунда при стандартных предпосылках и не переносятся автоматически на всю повторяющуюся систему с обучением.

В-четвёртых, задача адаптивного управления периферийными узлами формализована как децентрализованный частично наблюдаемый марковский процесс принятия решений. В предложенной постановке обучаемыми агентами являются периферийные вычислительные узлы, а не роботы, формирующие заявки. Действиями агентов являются операционные режимы узлов: режимы допуска, резервирования, приоритизации и защиты от перегрузки. Такое разделение ролей принципиально важно: обучаемая политика не управляет ставками и не изменяет оценки ценности задач, а значит, не подменяет аукционный механизм стратегическим обучением заявок.

В-пятых, обосновано применение кооперативного многоагентного обучения с подкреплением в парадигме централизованного обучения и децентрализованного исполнения. В качестве базового алгоритма выбран QMIX, поскольку он согласуется с дискретным пространством операционных режимов, командным вознаграждением, частичной наблюдаемостью и необходимостью локального исполнения политики периферийными узлами. Централизованное обучение может выполняться в серверном контуре, облаке, цифровом двойнике или имитационной среде, тогда как в режиме эксплуатации каждый узел выбирает режим на основе локального наблюдения.

В-шестых, разработан гибридный метод AMPLE-AMR. Метод объединяет обучаемую многоагентную политику выбора операционных режимов периферийных узлов и VCG-подобный аукционный слой распределения задач. В AMPLE-AMR обучаемая политика определяет условия распределения, то есть экспонируемую ёмкость и режимы обслуживания узлов, а аукционный слой выполняет назначение задач в рамках выбранных ресурсных ограничений. Такое построение позволяет совместить адаптацию к динамике среды и экономически интерпретируемое распределение дефицитных ресурсов.

В-седьмых, предложена масштабируемая версия C-AMPLE-AMR, основанная на кластеризации графа периферийной инфраструктуры и проведении локальных распределений внутри кластеров. Кластерная версия предназначена для случаев, когда глобальный аукционный раунд становится слишком дорогим по времени вычислений или объёму коммуникации. В работе введены критерии оценки целесообразности перехода к кластерному режиму, учитывающие накладные расходы глобального распределения, длительность управляющего слота, межкластерный срез графа и возможную потерю общественного благосостояния.

В-восьмых, реализован программный прототип AMPLE-AMR и сформирована воспроизводимая экспериментальная среда. Прототип включает генератор складских сценариев, модели роботов и периферийных узлов, эвристическое распределение, VCG-подобное распределение, кластерное VCG-подобное распределение, контур обучения режимов узлов и автоматический экспорт результатов в виде таблиц и графиков.

В-девятых, проведена экспериментальная проверка предложенного метода в сценариях стабильной нагрузки, пикового потока задач, разнородной периферийной инфраструктуры, деградации связи, отказов узлов, масштабирования и анализа чувствительности параметров режимов. Эксперименты показали, что наиболее устойчивый практический вклад вносит обучаемый выбор операционных режимов периферийных узлов. AMPLE-AMR превосходит Fixed-Heuristic по общественному благосостоянию во всех основных сценариях: на $124.0$ условных единицы при стабильной нагрузке, на $210.5$ при пиковой нагрузке, на $95.6$ в разнородной инфраструктуре, на $94.9$ при деградации связи и на $110.7$ при отказах узлов. При этом самостоятельное преимущество VCG-подобного распределения над жадной эвристикой при фиксированных режимах в текущей параметризации не подтверждено, а диагностическая роль внутренних цен подтверждена лишь частично.


Тем самым поставленная в диссертации цель достигнута: разработаны и исследованы модели и гибридный метод адаптивного распределения вычислительных задач автономных мобильных роботов между гетерогенными периферийными узлами. Предложенный метод сочетает VCG-подобное распределение задач и обучаемое управление операционными режимами периферийных узлов, сохраняя явное разделение между механизмом назначения задач и контуром адаптации.

Принципиально важно, что в предложенной архитектуре обучаемая компонента не подменяет механизм распределения, не обучает заявки роботов и не изменяет оценки ценности задач. Она отвечает только за выбор операционных режимов периферийных узлов и, следовательно, за адаптивную настройку условий распределения. Это позволяет сохранять экономическую интерпретируемость аукционного слоя и одновременно учитывать долговременную динамику очередей, задержек, отказов и нагрузки.

Экспериментальная проверка в воспроизводимой складской AMR-среде показала, что после разведения критериев сравнения самостоятельный вклад в повышение общественного благосостояния дают обе составляющие гибридного метода. При фиксированных режимах аукционный allocator Fixed-Auction превосходит latency-ориентированный baseline Fixed-Heuristic во всех 70 сопоставлениях со средним эффектом $+558.7$. Сам по себе обучаемый выбор режимов в QMIX-Heuristic даёт более умеренный, но устойчивый выигрыш $+82.2$, а полная гибридная схема AMPLE-AMR в среднем превосходит одно-компонентные варианты Fixed-Auction и QMIX-Heuristic на $+326.1$.

В стабильном сценарии общественное благосостояние возрастает от $634.4$ у Fixed-Heuristic до $1168.6$ у Fixed-Auction и до $1235.6$ у AMPLE-AMR. В сценарии пиковой нагрузки выигрыш ещё заметнее: $957.8 \rightarrow 2315.9 \rightarrow 2516.7$, причём AMPLE-AMR дополнительно уменьшает 95-й квантиль задержки относительно Fixed-Heuristic с $78.6$ до $66.4$~мс. В сценариях деградации связи и отказов узлов AMPLE-AMR также значительно повышает общественное благосостояние по сравнению с Fixed-Heuristic, соответственно с $100.0$ до $875.2$ и с $564.4$ до $1061.2$, одновременно улучшая хвост задержки. В сценарии разнородной инфраструктуры обнаруживается важный контрпример: наибольшее благосостояние достигается у Fixed-Auction ($1248.2$), тогда как AMPLE-AMR получает $1183.1$, что указывает на чувствительность обучаемой политики режимов к структуре гетерогенной периферийной среды.

Проверка гипотез показала, что H1, H2 и H3 подтверждаются: аукционный слой сам по себе повышает общественное благосостояние при фиксированных режимах, обучаемый выбор режимов также даёт положительный эффект, а полная гибридная схема в среднем превосходит одно-компонентные варианты. H4, H5 и H6 подтверждаются частично: внутренние цены оказываются лишь слабым диагностическим сигналом, а в пиковых и деградационных сценариях улучшение благосостояния и хвоста задержки сопровождается ухудшением completion rate и ростом drop rate по сравнению с throughput-ориентированной эвристикой. H7 не подтверждается: кластерная версия не является универсальной заменой глобальному распределению и оправдана только в наиболее крупных конфигурациях.

Такой результат является содержательно важным. Он показывает, что аукционный слой в текущей постановке следует рассматривать не только как интерпретируемый механизм выбора распределения, расчёта внутренних цен и анализа внешних эффектов, но и как самостоятельный оптимизатор общественного благосостояния. Одновременно эксперименты выявили содержательный компромисс между целями: эвристический baseline чаще обслуживает больше задач и удерживает меньший drop rate, тогда как аукционные методы выбирают более ценное подмножество задач и тем самым выигрывают по общественному благосостоянию. Следовательно, вклад обучаемой компоненты состоит не в подмене аукционного механизма, а в дополнительной адаптации условий распределения к динамике нагрузки, сетевых деградаций и отказов.

Отдельный практический вывод связан с масштабируемостью. Кластерная версия C-AMPLE-AMR не является универсальной заменой глобальному распределению: на малых и средних масштабах кластеризация чаще увеличивает накладные расходы. Однако на масштабе Warehouse-XL C-AMPLE-AMR снижает суммарные накладные расходы принятия решения с $264.2$ до $51.5$~мс, то есть примерно на $80.5\%$, при потере общественного благосостояния $41.9$. Следовательно, переход к кластерному режиму целесообразен только тогда, когда глобальный аукционный раунд начинает занимать значимую долю управляющего слота и становится потенциальным узким местом цикла управления.

Существенным итогом работы является также создание воспроизводимого программного прототипа и набора экспериментальных сценариев, позволяющих автоматически получать численные результаты, графики и \LaTeX-таблицы для анализа метода. Это делает предложенный подход проверяемым, расширяемым и пригодным для последующих сравнительных исследований, в которых могут использоваться новые аукционные механизмы, иные алгоритмы многоагентного обучения и более детализированные цифровые двойники роботизированной инфраструктуры.

Вместе с тем предложенный подход имеет ряд ограничений. Во-первых, строгие свойства VCG-подобного механизма относятся только к отдельному раунду распределения при фиксированных режимах узлов. Во-вторых, точное вычисление платежей и решение задачи выбора победителей могут быть дорогостоящими для крупных систем. В-третьих, качество обучаемой политики зависит от репрезентативности обучающих сценариев и параметризации имитационной среды. В-четвёртых, кластеризация снижает накладные расходы только при достаточно крупном масштабе и может приводить к потере качества распределения при существенных межкластерных связях. В-пятых, в текущей версии работы не рассматривается отдельный алгоритм выбора оркестратора распределения задач; этот вопрос вынесен за рамки основного метода.

Дальнейшие исследования могут быть направлены на:
\begin{enumerate}
	\item разработку алгоритма выбора оркестратора распределения задач, применимого как в глобальном, так и в кластерном режиме;
	\item исследование приближённых методов вычисления VCG-подобных платежей для крупных систем;
	\item дальнейшее исследование компромисса между максимизацией общественного благосостояния и метриками throughput/latency, включая более сложные ограничения, нелинейные стоимости и конфликтующие классы задач;
	\item расширение модели на случай делимых и распараллеливаемых задач;
	\item учёт неопределённости в оценках ресурсных требований и сроков обработки задач;
	\item перенос обученной политики между различными конфигурациями AMR-систем;
	\item проверку метода в более реалистичных робототехнических симуляторах и на данных, полученных из реальных систем управления роботами;
	\item исследование альтернативных алгоритмов многоагентного обучения, включая более выразительные методы факторизации функции ценности и актор-критик подходы;
	\item доработку кластерного распределения, включая более эффективный локальный аукцион и критерии динамического перехода между глобальным и кластерным режимом;
	\item разработку механизмов доверия, аудита и проверки корректности заявляемых параметров задач и ресурсов.
\end{enumerate}

Таким образом, диссертационная работа формирует основу для построения адаптивных, интерпретируемых и масштабируемых методов распределения вычислительных задач в системах автономных мобильных роботов с поддержкой периферийных вычислений. Предложенный метод AMPLE-AMR и его кластерная версия C-AMPLE-AMR объединяют аукционное распределение и кооперативное многоагентное обучение, сохраняя при этом явное разделение между экономическим механизмом назначения задач и обучаемым управлением режимами периферийных узлов.

В заключение автор выражает благодарность и большую признательность научному руководителю Корхову~В.\,В. за поддержку, помощь, обсуждение результатов и~научное руководство. Автор также благодарит всех, кто сделал настоящую работу возможной.

% Команда, неявное цитирование всех имеющихся в базах bib источников
\nocite{*}
